\documentclass[12pt]{article} %V5.1 - Corrected z Interpretation
\usepackage{amsfonts,amssymb,amsmath,mathrsfs,amsthm}
\usepackage{graphicx,caption}
\usepackage{epsfig,epstopdf}
\usepackage{nccmath}
\usepackage{float}
\usepackage[bottom=1in,top=1in,left=1in,right=1in]{geometry}

\usepackage[dvipsnames]{xcolor}
\usepackage{enumitem}
\usepackage{bbm}
\usepackage{mathtools}
\usepackage[utf8]{inputenc}
\usepackage[english]{babel}
\usepackage{setspace}
\onehalfspacing

\newtheorem{theorem}{Theorem}
\newtheorem{lemma}{Lemma}
\newtheorem{proposition}{Proposition}
\newtheorem{corollary}{Corollary}
\theoremstyle{remark}
\newtheorem{remark}{Remark}

\usepackage[breaklinks=true]{hyperref}
\definecolor{darkblue}{rgb}{0,0,0.25}
\hypersetup{
  colorlinks=true,
  citecolor=blue,
  linkcolor=darkblue,
  urlcolor=darkblue,
  pdfpagemode=none
}

\begin{document}

\noindent {\large Capital Holdup, Skill Supply, and Industrial Policy under Search Frictions}

\vspace{0.6cm}

\section*{GE labor wedge induced by the capital wedge}\label{app:GE_L_wedge}

\subsection*{Setup}\label{app:GE_L_setup}
Assume the Hosios condition holds, $\beta=\eta$. The equilibrium labor allocation $L_s$ is pinned down by the sectoral indifference condition
\begin{equation}\label{eq:SC_exact}
F(L_s, k_m, k_s) \;\equiv\; r(\epsilon_m - \epsilon_s) - \big[R_m(L_s, k_m, k_s) - R_s(L_s, k_m, k_s)\big] = 0.
\end{equation}
Worker returns are
\begin{equation}\label{eq:Ri_def}
R_i \;\equiv\; \beta \underbrace{\frac{\theta_i q_i(\theta_i)}{D_i(\theta_i)}}_{H_i(\theta_i)} \Pi_i,
\qquad 
\Pi_i \equiv \frac{p_i}{P}f_i(k_i) - z,
\end{equation}
where $z$ is the flow value of unemployment. The discount term is
\begin{equation}\label{eq:Di_def}
D_i(\theta_i)= r+s_i + (1-\beta_i)q_i(\theta_i)+\beta_i\theta_i q_i(\theta_i).
\end{equation}
Relative prices follow from the CES aggregator (with $\rho=1-1/\sigma$):
\begin{equation}\label{eq:CES_price}
\frac{p_i}{P} = \left( \frac{Y}{a_i^{1/\rho} Y_i} \right)^{1-\rho}.
\end{equation}
Market tightness is characterized by the sectoral tightness condition
\begin{equation}\label{eq:tightness_cond}
\frac{P_k k_i}{P}
= \frac{q_i(\theta_i)(1-\beta_i)}{r+s_i}\;
  \frac{\Pi_i}{D_i(\theta_i)}, \qquad i\in\{m,s\}.
\end{equation}

\subsection*{Exogenous Assumptions}\label{app:GE_L_assumptions}
We impose standard regularity conditions.

\begin{enumerate}[label=\textbf{(A\arabic*)}, leftmargin=2cm]
\item \label{ass:DMP} \textbf{Search:} $q_i(\theta)>0$ and $q_i'(\theta)<0$ for all relevant $\theta$, and $r+s_i>0$.
\item \label{ass:CES} \textbf{Production and CES:} $f_i'(k)>0$, $f_i''(k)<0$. The final good aggregator has elasticity of substitution $\sigma\in(0,\infty)$, so $\rho<1$.
\end{enumerate}

\subsection*{A.3 Lemmas}

\begin{lemma}[CES Derivatives]\label{lem:CES_derivs}
Let $\omega_i \equiv \frac{p_i Y_i}{P Y}$ denote the expenditure share. Under \eqref{eq:CES_price}, the elasticities of real prices with respect to sectoral outputs satisfy
\[
\frac{\partial \ln (p_i/P)}{\partial \ln Y_j} = (1-\rho)(\omega_j - \delta_{ij}),
\]
where $\delta_{ij}$ is the Kronecker delta ($\delta_{ij}=1$ if $i=j$ and $0$ otherwise).
\end{lemma}

\begin{proof}
Taking logs of \eqref{eq:CES_price} gives $\ln(p_i/P)=(1-\rho)(\ln Y-\ln Y_i)+\text{const.}$ Differentiating with respect to $\ln Y_j$ yields
\[
\frac{\partial \ln (p_i/P)}{\partial \ln Y_j}
=(1-\rho)\left(\frac{\partial \ln Y}{\partial \ln Y_j}-\delta_{ij}\right).
\]
Finally, $\partial \ln Y/\partial \ln Y_j=\omega_j$ follows from $\omega_j=(p_j/P)(Y_j/Y)$ and the final-good first-order condition $p_j/P=\partial Y/\partial Y_j$, which implies $\omega_j=(Y_j/Y)\,\partial Y/\partial Y_j=\partial \ln Y/\partial \ln Y_j$.
\end{proof}

\begin{lemma}[Return Monotonicity under the Tightness Condition]\label{lem:Psi_pos}
Under \eqref{eq:tightness_cond} and \textbf{(A1)}, worker return is strictly increasing in match surplus:
\[
\Psi_i \;\equiv\; \frac{dR_i}{d\Pi_i} \;>\; 0.
\]
\end{lemma}

\begin{proof}
By the chain rule, $\Psi_i=(dR_i/d\theta_i)/(d\Pi_i/d\theta_i)$. The tightness condition \eqref{eq:tightness_cond} implies
\[
\frac{\Pi_i}{D_i(\theta_i)}
=\frac{P_k k_i}{P}\cdot\frac{r+s_i}{(1-\beta_i)q_i(\theta_i)}.
\]
Substituting into \eqref{eq:Ri_def} gives
\[
R_i=\beta\,\theta_i q_i(\theta_i)\cdot\frac{\Pi_i}{D_i(\theta_i)}
=\beta\,\frac{P_k k_i}{P}\cdot\frac{r+s_i}{1-\beta_i}\,\theta_i,
\]
so $dR_i/d\theta_i>0$. Rearranging \eqref{eq:tightness_cond} also yields
\[
\Pi_i(\theta_i)
=\frac{P_k k_i}{P}\cdot\frac{r+s_i}{1-\beta_i}\,\frac{D_i(\theta_i)}{q_i(\theta_i)}.
\]
Differentiating and using \eqref{eq:Di_def} implies
\[
\frac{d\Pi_i}{d\theta_i}
=\frac{P_k k_i}{P}\cdot\frac{r+s_i}{1-\beta_i}\,
\frac{\beta_i q_i(\theta_i)^2-(r+s_i)q_i'(\theta_i)}{q_i(\theta_i)^2}>0,
\]
since $q_i(\theta_i)>0$, $q_i'(\theta_i)<0$, and $r+s_i>0$. Hence $\Psi_i>0$.
\end{proof}

\subsection*{A.4 Derivation of the Labor Response}\label{app:GE_L_deriv}

Applying the Implicit Function Theorem to \eqref{eq:SC_exact} yields
\begin{equation}\label{eq:IFT_master}
\frac{dL_s}{dk_m}
= - \frac{\partial F / \partial k_m}{\partial F / \partial L_s}
= - \frac{\frac{\partial (R_m - R_s)}{\partial k_m}}{\frac{\partial (R_m - R_s)}{\partial L_s}}.
\end{equation}

\noindent\textbf{1. The Stability Term (Denominator).}
We show that
\begin{equation}\label{eq:denom_goal}
\frac{\partial (R_m-R_s)}{\partial L_s} > 0 .
\end{equation}
Differentiating $R_i=\beta H_i(\theta_i)\Pi_i$ with respect to $L_s$ (holding $(k_m,k_s)$ fixed) gives
\begin{equation}\label{eq:dRi_dLs_explicit}
\frac{\partial R_i}{\partial L_s}
=\beta\left[\Pi_i H_i'(\theta_i)\frac{\partial \theta_i}{\partial L_s}
+H_i(\theta_i)\frac{\partial \Pi_i}{\partial L_s}\right].
\end{equation}
Since $\Pi_i=\frac{p_i}{P}f_i(k_i)-z$, holding $(k_i,z)$ fixed implies
\begin{equation}\label{eq:dPi_dLs}
\frac{\partial \Pi_i}{\partial L_s}=f_i(k_i)\frac{\partial (p_i/P)}{\partial L_s}.
\end{equation}
The tightness condition \eqref{eq:tightness_cond} induces an equilibrium mapping $\theta_i=\Theta_i(\Pi_i)$ for fixed $(k_i,P_k/P)$, so \eqref{eq:dRi_dLs_explicit} can be written along this locus as
\begin{equation}\label{eq:dRi_dLs_packaged}
\frac{\partial R_i}{\partial L_s}
=\frac{dR_i}{d\Pi_i}\frac{\partial \Pi_i}{\partial L_s}
\equiv \Psi_i\,\frac{\partial \Pi_i}{\partial L_s}.
\end{equation}
By Lemma \ref{lem:Psi_pos}, $\Psi_i>0$.

Under the resource constraint $L_m=\bar L-L_s$, an increase in $L_s$ raises $Y_s$ and lowers $Y_m$. Lemma \ref{lem:CES_derivs} therefore implies
\[
\frac{\partial (p_s/P)}{\partial L_s}<0,
\qquad
\frac{\partial (p_m/P)}{\partial L_s}>0,
\]
and hence, by \eqref{eq:dPi_dLs},
\[
\frac{\partial \Pi_s}{\partial L_s}<0,
\qquad
\frac{\partial \Pi_m}{\partial L_s}>0.
\]
Combining these inequalities with \eqref{eq:dRi_dLs_packaged} and $\Psi_i>0$ yields
\[
\frac{\partial R_s}{\partial L_s}<0,
\qquad
\frac{\partial R_m}{\partial L_s}>0,
\]
which in turn implies \eqref{eq:denom_goal}:
\[
\frac{\partial (R_m-R_s)}{\partial L_s}
=\frac{\partial R_m}{\partial L_s}-\frac{\partial R_s}{\partial L_s}
>0.
\]

\vspace{0.3cm}

% ============================================================
% A.4.2 Numerator and analytical condition for positivity
% (drop this directly after the denominator subsection)
% ============================================================

\subsubsection*{A.4.2 The capital-response term (numerator)}
Holding $(L_s,k_s)$ fixed, the numerator in \eqref{eq:IFT_master} is
\begin{equation}\label{eq:numer_def}
\frac{\partial (R_m-R_s)}{\partial k_m}
=\frac{\partial R_m}{\partial k_m}-\frac{\partial R_s}{\partial k_m}.
\end{equation}
Using \eqref{eq:Ri_def}, $R_i=\beta H_i(\theta_i)\Pi_i$ with
$H_i(\theta_i)=\theta_i q_i(\theta_i)/D_i(\theta_i)$ and
$\Pi_i=(p_i/P)f_i(k_i)-z$, the product and chain rules imply, for each
$i\in\{m,s\}$,
\begin{equation}\label{eq:dRi_dkm_general}
\frac{\partial R_i}{\partial k_m}
=\beta\left[
H_i(\theta_i)\frac{\partial \Pi_i}{\partial k_m}
+\Pi_i\,H_i'(\theta_i)\frac{\partial \theta_i}{\partial k_m}
\right].
\end{equation}
Substituting \eqref{eq:dRi_dkm_general} into \eqref{eq:numer_def} yields
\begin{equation}\label{eq:numer_expanded}
\frac{\partial (R_m-R_s)}{\partial k_m}
=\beta\Bigg(
H_m\,\frac{\partial \Pi_m}{\partial k_m}
+\Pi_m H_m'\,\frac{\partial \theta_m}{\partial k_m}
-
H_s\,\frac{\partial \Pi_s}{\partial k_m}
-\Pi_s H_s'\,\frac{\partial \theta_s}{\partial k_m}
\Bigg).
\end{equation}

\textbf{Derivation of $\frac{\partial \Pi_m}{\partial k_m}$}

Recall
\[
\Pi_i \equiv \frac{p_i}{P}f_i(k_i)-z,\qquad i\in\{m,s\}.
\]
Holding $(L_s,k_s)$ fixed, we have
\begin{align}
\frac{\partial \Pi_m}{\partial k_m}
&=
\frac{p_m}{P}\,f_m'(k_m)
+f_m(k_m)\,\frac{\partial}{\partial k_m}\!\left(\frac{p_m}{P}\right),
\label{eq:dPim_dkm}\\[0.4em]
\frac{\partial \Pi_s}{\partial k_m}
&=
f_s(k_s)\,\frac{\partial}{\partial k_m}\!\left(\frac{p_s}{P}\right),
\label{eq:dPis_dkm}
\end{align}
where the direct term in $f_s$ is zero because $k_s$ is held fixed.

\medskip

To express the price derivatives, use Lemma~\ref{lem:CES_derivs}:
\[
\frac{\partial \ln (p_i/P)}{\partial \ln Y_j}=(1-\rho)(\omega_j-\delta_{ij}),
\qquad \omega_j\equiv \frac{p_jY_j}{PY}.
\]
By the chain rule,
\[
\frac{\partial}{\partial k_m}\!\left(\frac{p_i}{P}\right)
=\frac{p_i}{P}\sum_{j\in\{m,s\}}\frac{\partial \ln (p_i/P)}{\partial \ln Y_j}\,
\frac{1}{Y_j}\frac{\partial Y_j}{\partial k_m}.
\]
Specializing to two sectors and simplifying gives
\begin{align}
\frac{\partial}{\partial k_m}\!\left(\frac{p_m}{P}\right)
&=\frac{p_m}{P}(1-\rho)\Bigg[
-(1-\omega_m)\frac{1}{Y_m}\frac{\partial Y_m}{\partial k_m}
+(1-\omega_m)\frac{1}{Y_s}\frac{\partial Y_s}{\partial k_m}
\Bigg],
\label{eq:dpmoPk}\\[0.4em]
\frac{\partial}{\partial k_m}\!\left(\frac{p_s}{P}\right)
&=\frac{p_s}{P}(1-\rho)\omega_m\Bigg[
\frac{1}{Y_m}\frac{\partial Y_m}{\partial k_m}
-\frac{1}{Y_s}\frac{\partial Y_s}{\partial k_m}
\Bigg].
\label{eq:dpsdPk}
\end{align}

Substituting \eqref{eq:dpmoPk} into \eqref{eq:dPim_dkm} and \eqref{eq:dpsdPk} into \eqref{eq:dPis_dkm} yields
\begin{align}
\frac{\partial \Pi_m}{\partial k_m}
&=
\frac{p_m}{P}f_m'(k_m)
+\frac{p_m}{P}f_m(k_m)(1-\rho)\Bigg[
-(1-\omega_m)\frac{1}{Y_m}\frac{\partial Y_m}{\partial k_m}
+(1-\omega_m)\frac{1}{Y_s}\frac{\partial Y_s}{\partial k_m}
\Bigg],
\label{eq:dPim_dkm_closed}\\[0.6em]
\frac{\partial \Pi_s}{\partial k_m}
&=
\frac{p_s}{P}f_s(k_s)(1-\rho)\omega_m\Bigg[
\frac{1}{Y_m}\frac{\partial Y_m}{\partial k_m}
-\frac{1}{Y_s}\frac{\partial Y_s}{\partial k_m}
\Bigg].
\label{eq:dPis_dkm_closed}
\end{align}

\textbf{Derivation of $\frac{\partial \theta_i}{\partial k_m}$}
This one is very complicated...
\end{document}